% GENERATED FILE. Edit canonical lessons, then run npm run generate:books.
% canonical-source-hash: fnv1a64:eb91b406c51a2686
% canonical-lessons: TA-C491-potai, TA-C491-pukaipiti, TA-C491-tiyanam, TA-C491-munaku, TA-C491-aruvai-cikiccai, TA-R491-review-485-491

\chapter{Habits and Healing}
\label{ch:ta-habits-and-healing}
\hlchaptermodality{491}

\begin{chapteropening}
\textbf{By the end of this chapter:} \emph{I can talk about addiction, smoking, meditation, groaning and surgery.}

Complete the last lesson of chapter 491: I can talk about addiction, smoking, meditation, groaning and surgery.
\end{chapteropening}

\section[\texorpdfstring{\ta{போதை} (pōtai)}{போதை (pōtai)}]{\ta{போதை} (pōtai) — intoxication; addiction}

\label{lesson:TA-C491-potai}



\begin{quote}
Before the new one: say the Tamil for to shake off, then the Tamil for to ache.
\end{quote}

\subsection*{You'll want to know: \ta{போதை}}
\textbf{\ta{போதை}} — \emph{pōtai} — \textquotedblleft{}intoxication; addiction\textquotedblright{}.

\begin{quote}
\ta{போதைப்} \ta{பழக்கம்} \ta{குடும்பத்தைக்} \ta{கெடுக்கும்}.
\end{quote}

\begin{quote}
\emph{pōtaip paḻakkam kuṭumpattaik keṭukkum.}
\end{quote}

\begin{quote}
Addiction ruins families.
\end{quote}

\ta{போதைப்} \ta{பழக்கம்}, an addiction.

\begin{cousinweb}
The origin of \textbf{\ta{போதை}} (\emph{pōtai}) is not settled; we could not find it in DEDR or CDIAL.
\end{cousinweb}

\begin{grammarlens}[title={What you've built}]
One more word for describing what you lived and felt.
\end{grammarlens}

\subsection*{Your turn}
\begin{itemize}
\raggedright
  \item \emph{Say it:} \emph{pōtai}
  \item \emph{Say it:} \emph{pōtai}, once more
  \item \emph{Say it:} \emph{pōtaip paḻakkam kuṭumpattaik keṭukkum.} — Addiction ruins families.
  \item \emph{Recall:} say the Tamil for to shake off, then the Tamil for to ache, then say \emph{pōtai} again
\end{itemize}

\begin{tcolorbox}[breakable,colback=teal!4,colframe=teal!35!black,title={Before you move on}]
What does \ta{போதை} mean? (\textquotedblleft{}Intoxication; addiction\textquotedblright{}.) Say it once more.
\end{tcolorbox}
\section[\texorpdfstring{\ta{புகைபிடி} (pukaipiṭi)}{புகைபிடி (pukaipiṭi)}]{\ta{புகைபிடி} (pukaipiṭi) — to smoke tobacco}

\label{lesson:TA-C491-pukaipiti}



\begin{quote}
Before the new one: say the Tamil for to ache, then the Tamil for intoxication.
\end{quote}

\subsection*{You'll want to know: \ta{புகைபிடி}}
\textbf{\ta{புகைபிடி}} — \emph{pukaipiṭi} — \textquotedblleft{}to smoke tobacco\textquotedblright{}.

\begin{quote}
\ta{என்} \ta{மாமா} \ta{புகைபிடிப்பதை} \ta{நிறுத்தினார்}.
\end{quote}

\begin{quote}
\emph{eṉ māmā pukaipiṭippatai niṟuttiṉār.}
\end{quote}

\begin{quote}
My uncle stopped smoking.
\end{quote}

Past (I …): \ta{புகைபிடித்தேன்}, \emph{pukaipiṭittēṉ}.

\ta{புகை} \textquotedblleft{}smoke\textquotedblright{} + \ta{பிடி} \textquotedblleft{}to catch\textquotedblright{}.

\begin{cousinweb}
\textbf{\ta{புகைபிடி}} (\emph{pukaipiṭi}) is a native compound of \ta{புகை} \emph{pukai}, \textquotedblleft{}smoke\textquotedblright{} (DEDR 4240), and \ta{பிடி} \emph{piṭi}, \textquotedblleft{}to catch, to hold\textquotedblright{} (DEDR 4148).
\end{cousinweb}

\begin{grammarlens}[title={What you've built}]
One more word for describing what you lived and felt.
\end{grammarlens}

\subsection*{Your turn}
\begin{itemize}
\raggedright
  \item \emph{Say it:} \emph{pukaipiṭi}
  \item \emph{Say it:} \emph{pukaipiṭi}, once more
  \item \emph{Say it:} \emph{eṉ māmā pukaipiṭippatai niṟuttiṉār.} — My uncle stopped smoking.
  \item \emph{Recall:} say the Tamil for to ache, then the Tamil for intoxication, then say \emph{pukaipiṭi} again
\end{itemize}

\begin{tcolorbox}[breakable,colback=teal!4,colframe=teal!35!black,title={Before you move on}]
What does \ta{புகைபிடி} mean? (\textquotedblleft{}To smoke tobacco\textquotedblright{}.) Say it once more.
\end{tcolorbox}
\section[\texorpdfstring{\ta{தியானம்} (tiyāṉam)}{தியானம் (tiyāṉam)}]{\ta{தியானம்} (tiyāṉam) — meditation}

\label{lesson:TA-C491-tiyanam}



\begin{quote}
Before the new one: say the Tamil for intoxication, then the Tamil for to smoke tobacco.
\end{quote}

\subsection*{You'll want to know: \ta{தியானம்}}
\textbf{\ta{தியானம்}} — \emph{tiyāṉam} — \textquotedblleft{}meditation\textquotedblright{}.

\begin{quote}
\ta{தினமும்} \ta{காலையில்} \ta{அவர்} \ta{பத்து} \ta{நிமிடம்} \ta{தியானம்} \ta{செய்கிறார்}.
\end{quote}

\begin{quote}
\emph{tiṉamum kālaiyil avar pattu nimiṭam tiyāṉam ceykiṟār.}
\end{quote}

\begin{quote}
Every morning he meditates for ten minutes.
\end{quote}

\ta{தியானம்} \ta{செய்}, to meditate.

\begin{cousinweb}
\textbf{\ta{தியானம்}} (\emph{tiyāṉam}) is borrowed from Sanskrit \emph{dhyāna}, \textquotedblleft{}meditation\textquotedblright{} (CDIAL 6889).
\end{cousinweb}

\begin{grammarlens}[title={What you've built}]
One more word for describing what you lived and felt.
\end{grammarlens}

\subsection*{Your turn}
\begin{itemize}
\raggedright
  \item \emph{Say it:} \emph{tiyāṉam}
  \item \emph{Say it:} \emph{tiyāṉam}, once more
  \item \emph{Say it:} \emph{tiṉamum kālaiyil avar pattu nimiṭam tiyāṉam ceykiṟār.} — Every morning he meditates for ten minutes.
  \item \emph{Recall:} say the Tamil for intoxication, then the Tamil for to smoke tobacco, then say \emph{tiyāṉam} again
\end{itemize}

\begin{tcolorbox}[breakable,colback=teal!4,colframe=teal!35!black,title={Before you move on}]
What does \ta{தியானம்} mean? (\textquotedblleft{}Meditation\textquotedblright{}.) Say it once more.
\end{tcolorbox}
\section[\texorpdfstring{\ta{முனகு} (muṉaku)}{முனகு (muṉaku)}]{\ta{முனகு} (muṉaku) — to groan; to moan}

\label{lesson:TA-C491-munaku}



\begin{quote}
Before the new one: say the Tamil for to smoke tobacco, then the Tamil for meditation.
\end{quote}

\subsection*{You'll want to know: \ta{முனகு}}
\textbf{\ta{முனகு}} — \emph{muṉaku} — \textquotedblleft{}to groan; to moan\textquotedblright{}.

\begin{quote}
\ta{வலியால்} \ta{நோயாளி} \ta{இரவு} \ta{முழுவதும்} \ta{முனகினார்}.
\end{quote}

\begin{quote}
\emph{valiyāl nōyāḷi iravu muḻuvatum muṉakiṉār.}
\end{quote}

\begin{quote}
The patient groaned all night in pain.
\end{quote}

Past (I …): \ta{முனகினேன்}, \emph{muṉakiṉēṉ}.

A low sound, softer than \ta{அலறு} \textquotedblleft{}to scream\textquotedblright{}.

\begin{cousinweb}
\textbf{\ta{முனகு}} (\emph{muṉaku}) is native Dravidian (DEDR 4856), in the same entry as \ta{முணுமுணு} \textquotedblleft{}to mutter\textquotedblright{}.
\end{cousinweb}

\begin{grammarlens}[title={What you've built}]
One more word for describing what you lived and felt.
\end{grammarlens}

\subsection*{Your turn}
\begin{itemize}
\raggedright
  \item \emph{Say it:} \emph{muṉaku}
  \item \emph{Say it:} \emph{muṉaku}, once more
  \item \emph{Say it:} \emph{valiyāl nōyāḷi iravu muḻuvatum muṉakiṉār.} — The patient groaned all night in pain.
  \item \emph{Recall:} say the Tamil for to smoke tobacco, then the Tamil for meditation, then say \emph{muṉaku} again
\end{itemize}

\begin{tcolorbox}[breakable,colback=teal!4,colframe=teal!35!black,title={Before you move on}]
What does \ta{முனகு} mean? (\textquotedblleft{}To groan; to moan\textquotedblright{}.) Say it once more.
\end{tcolorbox}
\section[\texorpdfstring{\ta{அறுவை} \ta{சிகிச்சை} (aṟuvai cikiccai)}{அறுவை சிகிச்சை (aṟuvai cikiccai)}]{\ta{அறுவை} \ta{சிகிச்சை} (aṟuvai cikiccai) — surgery; an operation}

\label{lesson:TA-C491-aruvai-cikiccai}



\begin{quote}
Before the new one: say the Tamil for meditation, then the Tamil for to groan.
\end{quote}

\subsection*{You'll want to know: \ta{அறுவை} \ta{சிகிச்சை}}
\textbf{\ta{அறுவை} \ta{சிகிச்சை}} — \emph{aṟuvai cikiccai} — \textquotedblleft{}surgery; an operation\textquotedblright{}.

\begin{quote}
\ta{அவருக்கு} \ta{இதயத்தில்} \ta{அறுவை} \ta{சிகிச்சை} \ta{நடந்தது}.
\end{quote}

\begin{quote}
\emph{avarukku itayattil aṟuvai cikiccai naṭantatu.}
\end{quote}

\begin{quote}
He had heart surgery.
\end{quote}

\ta{அறுவை} \textquotedblleft{}cutting\textquotedblright{} + \ta{சிகிச்சை} \textquotedblleft{}treatment\textquotedblright{}.

\begin{cousinweb}
\textbf{\ta{அறுவை} \ta{சிகிச்சை}} (\emph{aṟuvai cikiccai}) is a phrase built from native \ta{அறு} \emph{aṟu}, \textquotedblleft{}to cut\textquotedblright{}, and \ta{சிகிச்சை}, borrowed from Sanskrit \emph{cikitsā} (Monier-Williams).
\end{cousinweb}

\begin{grammarlens}[title={What you've built}]
One more word for describing what you lived and felt.
\end{grammarlens}

\subsection*{Your turn}
\begin{itemize}
\raggedright
  \item \emph{Say it:} \emph{aṟuvai cikiccai}
  \item \emph{Say it:} \emph{aṟuvai cikiccai}, once more
  \item \emph{Say it:} \emph{avarukku itayattil aṟuvai cikiccai naṭantatu.} — He had heart surgery.
  \item \emph{Recall:} say the Tamil for meditation, then the Tamil for to groan, then say \emph{aṟuvai cikiccai} again
\end{itemize}

\begin{tcolorbox}[breakable,colback=teal!4,colframe=teal!35!black,title={Before you move on}]
What does \ta{அறுவை} \ta{சிகிச்சை} mean? (\textquotedblleft{}Surgery; an operation\textquotedblright{}.) Say it once more.
\end{tcolorbox}
\section[(dialogue)]{Looking back: from Home Remedies to Habits and Healing}

\label{lesson:TA-R491-review-485-491}



\begin{quote}
Say the Tamil for a herbal decoction and dry ginger.
\end{quote}

\subsection*{Your turn}
Say these aloud:

\begin{itemize}
\raggedright
  \item a herbal decoction, dry ginger, rice gruel, a medicated oil, an ointment
  \item to rub on, to press down, to gargle, a walking stick, to pull through
  \item body fat, protein, a vitamin, to digest, to burp
  \item fatness, a paunch, to grow thin, a hunch, starvation
  \item to slumber, to snore, to toss and turn, to wake up, to perspire
  \item to straighten up, to bend down, to lean, to shake off, to ache
  \item intoxication, to smoke tobacco, meditation, to groan, surgery
\end{itemize}

\begin{tcolorbox}[breakable,colback=teal!4,colframe=teal!35!black,title={Before you move on}]
A herbal decoction? (\textbf{\ta{கஷாயம்}}, \emph{kaṣāyam}.) Dry ginger? (\textbf{\ta{சுக்கு}}, \emph{cukku}.) Rice gruel? (\textbf{\ta{கஞ்சி}}, \emph{kañci}.) A medicated oil? (\textbf{\ta{தைலம்}}, \emph{tailam}.) An ointment? (\textbf{\ta{களிம்பு}}, \emph{kaḷimpu}.) To rub on? (\textbf{\ta{தடவு}}, \emph{taṭavu}.) To press down? (\textbf{\ta{அமுக்கு}}, \emph{amukku}.) To gargle? (\textbf{\ta{கொப்பளி}}, \emph{koppaḷi}.) A walking stick? (\textbf{\ta{ஊன்றுகோல்}}, \emph{ūṉṟukōl}.) To pull through? (\textbf{\ta{மீள்}}, \emph{mīḷ}.) Body fat? (\textbf{\ta{கொழுப்பு}}, \emph{koḻuppu}.) Protein? (\textbf{\ta{புரதம்}}, \emph{puratam}.) A vitamin? (\textbf{\ta{வைட்டமின்}}, \emph{vaiṭṭamiṉ}.) To digest? (\textbf{\ta{செரி}}, \emph{ceri}.) To burp? (\textbf{\ta{ஏப்பம்} \ta{விடு}}, \emph{ēppam viṭu}.) Fatness? (\textbf{\ta{பருமன்}}, \emph{parumaṉ}.) A paunch? (\textbf{\ta{தொப்பை}}, \emph{toppai}.) To grow thin? (\textbf{\ta{மெலி}}, \emph{meli}.) A hunch? (\textbf{\ta{கூன்}}, \emph{kūṉ}.) Starvation? (\textbf{\ta{பட்டினி}}, \emph{paṭṭiṉi}.) To slumber? (\textbf{\ta{உறங்கு}}, \emph{uṟaṅku}.) To snore? (\textbf{\ta{குறட்டை} \ta{விடு}}, \emph{kuṟaṭṭai viṭu}.) To toss and turn? (\textbf{\ta{புரள்}}, \emph{puraḷ}.) To wake up? (\textbf{\ta{விழி}}, \emph{viḻi}.) To perspire? (\textbf{\ta{வியர்}}, \emph{viyar}.) To straighten up? (\textbf{\ta{நிமிர்}}, \emph{nimir}.) To bend down? (\textbf{\ta{குனி}}, \emph{kuṉi}.) To lean? (\textbf{\ta{சாய்}}, \emph{cāy}.) To shake off? (\textbf{\ta{உதறு}}, \emph{utaṟu}.) To ache? (\textbf{\ta{நோ}}, \emph{nō}.) Intoxication? (\textbf{\ta{போதை}}, \emph{pōtai}.) To smoke tobacco? (\textbf{\ta{புகைபிடி}}, \emph{pukaipiṭi}.) Meditation? (\textbf{\ta{தியானம்}}, \emph{tiyāṉam}.) To groan? (\textbf{\ta{முனகு}}, \emph{muṉaku}.) Surgery? (\textbf{\ta{அறுவை} \ta{சிகிச்சை}}, \emph{aṟuvai cikiccai}.)
\end{tcolorbox}
